\documentclass{article}
\usepackage[utf8]{inputenc}
\usepackage{amsmath}
\usepackage{amssymb}
\usepackage{geometry}
\usepackage{enumerate}

\geometry{a4paper, margin=1in}

\begin{document}

\section*{CSE400 – Fundamentals of Probability in Computing}
\subsection*{Lecture 21 and 22: Inequalities and Tail Bound}
\textit{(Prepared as an exam-oriented lecture scribe from the provided lecture slides)}

\subsection*{1. Topic Introduction}
The lecture is structured around the study of the following key concepts:
\begin{itemize}
    \item Tail and motivation of tail bounds
    \item Markov’s Inequality
    \item Chebyshev’s Inequality
    \item Chernoff Bound and Poisson Tail
    \item Jensen’s Inequality
    \item Convex Functions
    \item Case Study: System-level understanding
\end{itemize}

\subsection*{2. Logical Flow of Topics}
The sequence of topics as presented in the slides follows a specific progression:

\paragraph{Step 1: Tail and Motivation of Tail Bound}
\begin{itemize}
    \item Introduces the concept of tail behavior.
    \item Establishes the motivation for studying bounds on tail probabilities.
\end{itemize}

\paragraph{Step 2: Markov’s Inequality}
\begin{itemize}
    \item Includes Definition, Proof, and Example.
\end{itemize}

\paragraph{Step 3: Chebyshev’s Inequality}
\begin{itemize}
    \item Includes Definition, Proof, and Example.
\end{itemize}

\paragraph{Step 4: Chernoff Bound and Poisson Tail}
\begin{itemize}
    \item Includes Definition, Proof, and Example.
\end{itemize}

\paragraph{Step 5: Jensen’s Inequality}
\begin{itemize}
    \item Focuses on Convex functions, including Proof and Example.
\end{itemize}

\paragraph{Step 6: Case Study}
\begin{itemize}
    \item \textbf{Problem Statement:} How to predict the worst-case demands on a server?
    \item This indicates the application of tail bounds in system-level performance analysis.
\end{itemize}

\subsection*{3. Structural Reasoning Flow}
The lecture follows a consistent reasoning structure for each inequality introduced:
\begin{enumerate}
    \item Introduce the inequality.
    \item Provide its definition.
    \item Present the proof.
    \item Demonstrate via example.
\end{enumerate}
This pattern is applied to Markov’s Inequality, Chebyshev’s Inequality, and the Chernoff Bound, then extended to Jensen’s Inequality with convexity.

\subsection*{4. Closing}
The lecture concludes with:
\begin{itemize}
    \item A system-level case study.
    \item Application of inequalities to worst-case analysis.
    \item Final summary slide.
\end{itemize}

\subsection*{Important Note for Exam Preparation}
The slides specify the structure and topics but do not provide the internal derivations or formulas. The examinable components identifiable from the material are:
\begin{itemize}
    \item Names of inequalities.
    \item Their sequence and grouping.
    \item Association with definitions, proofs, and examples.
    \item Application to tail bounds and worst-case analysis.
\end{itemize}

\end{document}